\begingroup
\renewcommand{\arraystretch}{1.45}
\fancyhead[L]{\small ARITHMETIC–GEOMETRIC MEAN.}
\Needspace{30\baselineskip}
\AGMArticle{17}
[From the properties derived in the preceding article of the functions $p,q,r$ defined by the series stated there, it follows that, if $a,b,c$ are three quantities satisfying the equation $aa=bb+cc$, they can be put in the form
\Needspace{10\baselineskip}
\[
\begin{gathered}a=\mathrm M(a,b)\cdot(py)^2,\\[5pt] b=\mathrm M(a,b)\cdot(qy)^2,\\[5pt] c=\mathrm M(a,b)\cdot(ry)^2\end{gathered}
\]
and that then $\frac{\log y}{\pi}=-\frac{\mathrm M(a,b)}{\mathrm M(a,c)}$ must hold. If one further sets
\Needspace{10\baselineskip}
\[
\begin{gathered}a=\mathrm M(a,c)\cdot(pz)^2,\\[5pt] c=\mathrm M(a,c)\cdot(qz)^2,\\[5pt] b=\mathrm M(a,b)\cdot(rz)^2\end{gathered}
\]
then
\Needspace{10\baselineskip}
\[
\frac{\log z}{\pi}=-\frac{\mathrm M(a,c)}{\mathrm M(a,b)}=\frac{\pi}{\log y}.
\]
The theorem obtained by combining these two representations was stated by \textsc{Gauss} in the form that the functions
\Needspace{10\baselineskip}
\[
\begin{aligned}\mathfrak P t&=1+2e^{-\pi t}+2e^{-4\pi t}+2e^{-9\pi t}\\&\quad+.\;.+2e^{-nn\pi t}+.\;.\end{aligned}
\]
\Needspace{10\baselineskip}
\[
\begin{aligned}\mathfrak Q t&=1-2e^{-\pi t}+2e^{-4\pi t}-2e^{-9\pi t}\\&\quad+.\;.\pm2e^{-nn\pi t}\mp.\;.\end{aligned}
\]
\Needspace{10\baselineskip}
\[
\begin{aligned}\mathfrak R t&=2e^{-\frac14\pi t}+2e^{-\frac94\pi t}+2e^{-\frac{25}4\pi t}\\&\quad+.\;.+2e^{-(n+\frac12)^2\pi t}+.\;.\end{aligned}
\]
\Needspace{10\baselineskip}
satisfy the equations
\SourceLeaf{395}
\[
\begin{gathered}\mathfrak{P}t=\frac1{\surd t}\mathfrak{P}\frac1t,\\[5pt]
\mathfrak{Q}t=\frac1{\surd t}\mathfrak{R}\frac1t,\\[5pt]
\mathfrak{R}t=\frac1{\surd t}\mathfrak{Q}\frac1t\end{gathered}
\]
where the signs of the square roots are to be taken so that the real part is positive.

\Needspace{25\baselineskip}
From this property and the other property of the same functions that he recorded, namely that
\Needspace{10\baselineskip}
\[
\begin{gathered}\mathfrak{P}t=\mathfrak{Q}(t+i),\\[5pt]
\mathfrak{Q}t=\mathfrak{P}(t+i),\\[5pt]
\mathfrak{R}t=\surd i.\mathfrak{R}(t+i)\end{gathered}
\]
holds, \textsc{Gauss} seems to have derived the following theorem:]

\begin{itshape}
Let $\alpha,\allowbreak \beta,\allowbreak \gamma,\allowbreak \delta$ be real integers, $\alpha\delta-\beta\gamma=1$,
$\dfrac{\alpha t-\beta i}{\delta+\gamma ti}=t'$.

\Needspace{35\baselineskip}
We distinguish 6 cases, according as, modulo 2,
\Needspace{18\baselineskip}
\[
\begin{array}{r@{\;}c@{\quad}cccccc}
\alpha&\equiv&1&1&1&0&1&0\\
\beta&\equiv&0&1&0&1&1&1\\
\gamma&\equiv&0&0&1&1&1&1\\
\delta&\equiv&1&1&1&1&0&0
\end{array}
\]
Then
\Needspace{18\baselineskip}
\[
\begin{array}{r@{\;}c@{\quad}cccccc}
h\mathfrak{P}t'&=&\mathfrak{P}t&\mathfrak{Q}t&\mathfrak{R}t&\mathfrak{R}t&\mathfrak{Q}t&\mathfrak{P}t\\
h\mathfrak{Q}t'&=&\mathfrak{Q}t&\mathfrak{P}t&\mathfrak{Q}t&\mathfrak{P}t&\mathfrak{R}t&\mathfrak{R}t\\
h\mathfrak{R}t'&=&\mathfrak{R}t&\mathfrak{R}t&\mathfrak{P}t&\mathfrak{Q}t&\mathfrak{P}t&\mathfrak{Q}t
\end{array}
\]
\Needspace{10\baselineskip}
\[
h=\surd i^\lambda(\delta+\gamma ti)
\]
\end{itshape}

[where $\lambda$ generally has different values for the factors of the three functions $\mathfrak{P},\allowbreak \mathfrak{Q},\allowbreak \mathfrak{R}$.]

\begin{itshape}
\Needspace{24\baselineskip}
If here $t=\dfrac{\surd d+bi}{a}$, $t'=\dfrac{\surd d+b'i}{a'}$,
$-d=bb-ac=b'b'-a'c'$, then the form $(a,\allowbreak b,\allowbreak c)$ passes into $(a',\allowbreak b',\allowbreak c')$ through the transformation
$\left(\begin{smallmatrix}\delta,\allowbreak &-\beta\\-\gamma,\allowbreak &\alpha\end{smallmatrix}\right)$.

\Needspace{14\baselineskip}
Connection between the forms of the negative determinant $-p$ and the summatory functions.

Namely, if the forms $(a,\allowbreak b,\allowbreak c)$ $(A,\allowbreak B,\allowbreak C)$ are equivalent, the function $f$ is to be considered for which $ft\equiv fu$ both when $\dfrac{t-u}{i}$ is an integer and when
$t=\dfrac1u$.
\par
Each class then corresponds to a definite value of $f\dfrac{\surd p+bi}{a}$.
\end{itshape}

\SourceLeaf{396}
\Needspace{34\baselineskip}
\AGMArticle{18}

[\textsc{Gauss} connected another algorithm with that of the arithmetic–geometric mean. Like the latter, it starts from two given quantities, here to be denoted by $\alpha$ and $\beta$, and can be brought into a form showing many analogies with it, namely, if
\Needspace{10\baselineskip}
\[
\begin{aligned}
4a'a'&=(a+b)^2,\quad b'b'=ab,\quad\text{etc.}\\
4\alpha'\alpha'&=(\alpha+\beta)^2,\quad\beta'\beta'=\alpha\beta,\quad\text{etc.}
\end{aligned}
\]
is set. Just as the preceding investigations became considerably easier to survey by introducing, through the equation $aa=bb+cc$, the quantity $c$, a corresponding simplification of the formulas is obtained here if
$\gamma$ is introduced by the equation $\alpha\alpha=\beta\beta+\gamma\gamma$ and $\delta$ by
\Needspace{10\baselineskip}
\[
\begin{aligned}bc(b\gamma-c\beta)&=ca(a\gamma-c\alpha)\\&=ab(b\alpha-a\beta)\\&=abc\delta\end{aligned}
\]
and likewise $\gamma^{\mathrm n}$, $\delta^{\mathrm n}$ by the same equations after all symbols have been given the index $n$. Among the relations between these quantities, the following find frequent application in investigating this algorithm:
\Needspace{26\baselineskip}
\[
\begin{gathered}\frac{\alpha+\beta}{a+b}=\frac{\gamma-\delta}{c},\\[5pt]
\frac{\alpha-\beta}{a-b}=\frac{\gamma+\delta}{c},\\[5pt]
\frac{\alpha+\gamma}{a+c}=\frac{\beta+\delta}{b},\\[5pt]
\frac{\alpha-\gamma}{a-c}=\frac{\beta-\delta}{c},\end{gathered}
\]
\Needspace{26\baselineskip}
\[
\begin{aligned}
\alpha'&=\frac1{a'}\left(\frac{\alpha+\beta}{2}\right)^2
=\frac1{c'}\left(\frac{\gamma-\delta}{2}\right)^2,\quad
\beta'=\frac1{b'}\alpha\beta\\
\gamma'&=\frac1{c'}\left(\frac{\alpha-\beta}{2}\right)^2
=\frac1{a'}\left(\frac{\gamma+\delta}{2}\right)^2,\quad
\delta'=\frac1{b'}\gamma\delta
\end{aligned}
\]
\Needspace{18\baselineskip}
\[
\begin{aligned}\alpha\alpha+\delta\delta&=\beta\beta+\gamma\gamma
\\&=\frac ab(\alpha\beta+\gamma\delta)\\&=\frac ac(\alpha\gamma-\beta\delta)
\\&=a(\alpha'+\gamma')\end{aligned}
\]
\Needspace{18\baselineskip}
\[
\begin{aligned}\alpha\alpha-\gamma\gamma&=\beta\beta-\delta\delta
\\&=\frac bc(\beta\gamma-\alpha\delta)\\&=\frac ba(\alpha\beta-\gamma\delta)
\\&=b(\alpha'-\gamma')\end{aligned}
\]
\Needspace{18\baselineskip}
\[
\begin{aligned}\alpha\alpha-\beta\beta&=\gamma\gamma-\delta\delta
\\&=\frac ca(\alpha\gamma+\beta\delta)\\&=\frac cb(\beta\gamma+\alpha\delta)
\\&=2c\surd{\alpha'\gamma'}\end{aligned}
\]
\Needspace{18\baselineskip}
The limits of the terms in the seven progressions of quantities $a,\allowbreak b,\allowbreak c,\allowbreak \alpha,\allowbreak \beta,\allowbreak \gamma,\allowbreak \delta$ can be reduced to these four:

\SourceLeaf{397}
\Needspace{10\baselineskip}
\[
k=\mathrm M(a,b)=\lim a^{\mathrm n}=\lim b^{\mathrm n}
\]
\Needspace{26\baselineskip}
\[
\begin{aligned}
\surd y&=e^{-\frac\pi2\frac{\mathrm M(a,b)}{\mathrm M(a,c)}}
=\lim{}^{2^n}\!\surd{\frac{c^{\mathrm n}}{4a^{\mathrm n}}}\\
&=\frac{c}{4a}\cdot\frac a{a'}\cdot\surd{\frac{a'}{a''}}
\cdot{}^{2^2}\!\surd{\frac{a''}{a'''}}\\&\quad\cdot{}^{2^3}\!\surd{\frac{a'''}{a''''}}
.\;.{}^{2^{n-1}}\!\surd{\frac{a^{\mathrm n-1}}{a^{\mathrm n}}}.\;.
\end{aligned}
\]
\Needspace{26\baselineskip}
\[
\begin{aligned}
\frac xk&=\lim{}^{2^n}\!\surd{\frac{\alpha^{\mathrm n}}{k}}
=\lim{}^{2^n}\!\surd{\frac{\beta^{\mathrm n}}{k}}\\
&=\frac\beta b\cdot\surd{\frac ba\frac\alpha\beta}
\cdot{}^{2^2}\!\surd{\frac{b'}{a'}\frac{\alpha'}{\beta'}}
\\&\quad\cdot{}^{2^3}\!\surd{\frac{b''}{a''}\frac{\alpha''}{\beta''}}
.\;.\;.{}^{2^n}\!\surd{\frac{b^{\mathrm n-1}}{a^{\mathrm n-1}}\frac{\alpha^{\mathrm n-1}}{\beta^{\mathrm n-1}}}.\;.
\end{aligned}
\]
\Needspace{10\baselineskip}
\[
\frac xk\surd y.\eta^{\pm1}
=\lim{}^{2^{n-1}}\!\surd{\left[\surd{\frac{\gamma^{\mathrm n}}{4k}}\pm\surd{\frac{\delta^{\mathrm n}}{4k}}\right]}
\]
\Needspace{18\baselineskip}
If $\alpha$ and $\beta$ and all quantities $a^{\mathrm n}$, $b^{\mathrm n}$ are positive, then
$\dfrac{a^{\mathrm n}}{b^{\mathrm n}}$ and $\dfrac{a^{\mathrm n}a^{\mathrm n}}{b^{\mathrm n}\beta^{\mathrm n}}$ decrease steadily from the values
$\dfrac{a'}{b'}$ and $\dfrac{a'a'}{b'\beta'}$ to unity;
the foregoing expressions therefore give a definite value for $x$.
The same follows for $\eta$, if $\gamma$, $\delta$ and all $c^{\mathrm n}$ and $b^{\mathrm n}$ are positive, from
\Needspace{26\baselineskip}
\[
\begin{aligned}
\frac xk y^{\frac12}\eta&=\lim{}^{2^n}\!\surd{\frac{\gamma^{\mathrm n}}{k}}
=\lim{}^{2^n}\!\surd{\frac{\delta^{\mathrm n}}{k}}\\
&=\frac\delta b\cdot\surd{\frac ba\frac\gamma\delta}
\cdot{}^{2^2}\!\surd{\frac{b'}{a'}\frac{\gamma'}{\delta'}}
\\&\quad\cdot{}^{2^3}\!\surd{\frac{b''}{a''}\frac{\gamma''}{\delta''}}
.\;.\;.{}^{2^n}\!\surd{\frac{b^{\mathrm n-1}}{a^{\mathrm n-1}}\frac{\gamma^{\mathrm n-1}}{\delta^{\mathrm n-1}}}.\;.\;.
\end{aligned}
\]
\Needspace{25\baselineskip}
but, if $\delta$ is negative, from the substitution used by \textsc{Gauss}
\Needspace{26\baselineskip}
\[
\begin{gathered}
\surd{\frac{-\delta}{\gamma}}
=\mathrm{tang}\,U\cdot\surd{\frac ab}
=\mathrm{tang}\,V\cdot\surd{\frac ba}
\\=\surd{(\mathrm{tang}\,U\cdot\mathrm{tang}\,V)}
=\mathrm{tang}\,U'\\
U+V=2V'
\end{gathered}
\]
\Needspace{26\baselineskip}
\[
\begin{gathered}
\surd{\frac{-\delta'}{\gamma'}}
=\mathrm{tang}\,2U'\cdot\surd{\frac{a'}{b'}}
=\mathrm{tang}\,2V'\cdot\surd{\frac{b'}{a'}}
\\=\surd{(\mathrm{tang}\,2U'\cdot\mathrm{tang}\,2V')}
=\mathrm{tang}\,2U''\\
U'+V'=2V''
\end{gathered}
\]
\Needspace{26\baselineskip}
\[
\begin{gathered}
\surd{\frac{-\delta''}{\gamma''}}
=\mathrm{tang}\,4U''\cdot\surd{\frac{a''}{b''}}
=\mathrm{tang}\,4V''\cdot\surd{\frac{b''}{a''}}
\\=\surd{(\mathrm{tang}\,4U''\cdot\mathrm{tang}\,4V'')}
=\mathrm{tang}\,4U'''\\
U''+V''=2V'''
\end{gathered}
\]
\Needspace{26\baselineskip}
\[
\begin{gathered}
\surd{\frac{-\delta^{\mathrm n}}{\gamma^{\mathrm n}}}
=\mathrm{tang}\,2^nU^{\mathrm n}\cdot\surd{\frac{a^{\mathrm n}}{b^{\mathrm n}}}
=\mathrm{tang}\,2^nV^{\mathrm n}\cdot\surd{\frac{b^{\mathrm n}}{a^{\mathrm n}}}
\\=\surd{(\mathrm{tang}\,2^nU^{\mathrm n}\cdot\mathrm{tang}\,2^nV^{\mathrm n})}
=\mathrm{tang}\,2^nU^{\mathrm n+1}\\
U^{\mathrm n}+V^{\mathrm n}=2V^{\mathrm n+1}
\end{gathered}
\]
because then
\Needspace{32\baselineskip}
\[
\begin{aligned}
\sin U^2&=\frac bc\frac{-\delta}{\alpha},\\[4pt]&\quad
\cos U^2=\frac ac\frac\gamma\alpha,\\[4pt]&\quad
aa\sin U^2+bb\cos U^2=ab\frac\beta\alpha\\[4pt]
\sin V^2&=\frac ac\frac{-\delta}{\beta},\\[4pt]&\quad
\cos V^2=\frac bc\frac\gamma\beta,\\[4pt]&\quad
aa\cos V^2+bb\sin V^2=ab\frac\alpha\beta
\end{aligned}
\]
holds, and these equations also hold when $a^{\mathrm n},\allowbreak b^{\mathrm n},\allowbreak c^{\mathrm n},\allowbreak \alpha^{\mathrm n},\allowbreak \beta^{\mathrm n},\allowbreak \gamma^{\mathrm n},\allowbreak \delta^{\mathrm n},\allowbreak 2^nU^{\mathrm n},\allowbreak 2^nV^{\mathrm n}$
are substituted for $a,\allowbreak b,\allowbreak c,\allowbreak \alpha,\allowbreak \beta,\allowbreak \gamma,\allowbreak \delta,\allowbreak U,\allowbreak V$, so that
\Needspace{10\baselineskip}
\[
\eta^{\pm\frac12}
=\lim e^{\pm iU^{\mathrm n}}=\lim e^{\pm iV^{\mathrm n}}=e^{\pm iu}
\]
results.

\SourceLeaf{398}
\Needspace{24\baselineskip}
For $\delta=0$ all $\delta^{\mathrm n}$ vanish, and
\Needspace{26\baselineskip}
\[
\begin{gathered}\frac\alpha a=\frac\beta b=\frac\gamma c=.\;.
\\={}^{2^n}\!\surd{\frac{\alpha^{\mathrm n}}{a^{\mathrm n}}}
={}^{2^n}\!\surd{\frac{\beta^{\mathrm n}}{b^{\mathrm n}}}
={}^{2^n}\!\surd{\frac{\gamma^{\mathrm n}}{c^{\mathrm n}}}=\frac xk,\\[5pt]
{}^{2^n}\!\surd{\frac{\gamma^{\mathrm n}}{\alpha^{\mathrm n}}\frac{a^{\mathrm n}}{c^{\mathrm n}}}=\eta=1,\\[5pt] u=0\end{gathered}
\]
\Needspace{25\baselineskip}
If the limits $k,\allowbreak y,\allowbreak x,\allowbreak \eta,\allowbreak u$ obtained when the combined algorithm is formed starting from the quantities $a,\allowbreak b,\allowbreak \alpha,\allowbreak \beta$ are compared with the limits $k',\allowbreak y',\allowbreak x',\allowbreak \eta',\allowbreak u'$ that would be obtained by starting such an algorithm from the definite quantities $a',\allowbreak b',\allowbreak \alpha',\allowbreak \beta'$ previously obtained as initial terms, it is immediately seen that
\Needspace{18\baselineskip}
\[
\begin{gathered}k'=k,\\[5pt] y'=yy,\\[5pt]\frac{x'}{k'}=\frac{xx}{kk},\\[5pt]
\eta'=\eta\eta,\\[5pt] u'=2u\end{gathered}
\]
must hold, and that the equations formed by this law give, for $a^{\mathrm n},\allowbreak b^{\mathrm n},\allowbreak \alpha^{\mathrm n},\allowbreak \beta^{\mathrm n}$, the corresponding limits $k^{\mathrm n},\allowbreak y^{\mathrm n},\allowbreak x^{\mathrm n},\allowbreak \eta^{\mathrm n},\allowbreak u^{\mathrm n}$.

\Needspace{25\baselineskip}
From equations of the form:
\Needspace{18\baselineskip}
\[
\begin{gathered}\frac{c'}{a'}=\frac{a-b}{a+b},\\[5pt]
\frac{c'}{a'}\frac{\gamma'}{\alpha'}=\left(\frac{\alpha-\beta}{\alpha+\beta}\right)^2,\\[5pt]
\frac{c'}{a'}\frac{\alpha'}{\gamma'}=\left(\frac{\gamma-\delta}{\gamma+\delta}\right)^2\end{gathered}
\]
it follows that, if all the terms of the algorithm become positive, $y$, $y\eta$ and
$\dfrac1\eta$ are less than unity.

\Needspace{28\baselineskip}
The method given by \textsc{Gauss} for determining the limit of
$U^{\mathrm n}$ and $V^{\mathrm n}$, and likewise the determination of $\dfrac xk$ with the aid of those angles, leads very rapidly to the result in numerical calculations. The formulas above for determining a real $\eta$ and the corresponding $\dfrac xk$ are less convenient for numerical calculation. This circumstance will have been the reason why \textsc{Gauss} set out the algorithm:]
\Needspace{10\baselineskip}
\[
\begin{aligned}
A'&=\frac{A+B}{2},\quad B'=\frac{2ABa'}{b'(A+B)}\\
A''&=\frac{A'+B'}{2},\quad B''=\frac{2A'B'a''}{b'(A'+B')}
\end{aligned}
\]
and so on. [with the limit]
\Needspace{10\baselineskip}
\[
\begin{aligned}\frac Hk&=\frac Bb\surd{\frac{bA}{aB}}
\cdot{}^{4}\!\surd{\frac{b'A'}{a'B'}}
\\&\quad\cdot{}^{8}\!\surd{\frac{b''A''}{a''B''}}
\cdot{}^{16}\!\surd{\frac{b'''A'''}{a'''B'''}}.\;.\end{aligned}
\]
\Needspace{25\baselineskip}
[which can be reduced to the one above if
$A=\alpha$, $B=\beta$ are set, since then
\SourceLeaf{399}
\Needspace{26\baselineskip}
\[
\begin{gathered}A^{\mathrm n}=\dfrac{b^{\mathrm n+1}}{b'}\surd{\dfrac{\alpha^{\mathrm n}}{\beta^{\mathrm n}}}\cdot\surd{\alpha\beta},\\[9pt]
B^{\mathrm n}=\dfrac{b^{\mathrm n+1}}{b'}\surd{\dfrac{\beta_{\mathrm n}}{\alpha^{\mathrm n}}}\cdot\surd{\alpha\beta},\\[9pt] H=x\end{gathered}
\]
holds. The determination of a real $\eta$ is left unfinished by a gap in \textsc{Gauss}'s notes. It is obtained, however, if the same algorithm is formed from $C=\gamma$ and $D=\delta$ as just from $A$ and $B$; for then:
\Needspace{26\baselineskip}
\[
\begin{gathered}
C^{\mathrm n}=\dfrac{b^{\mathrm n+1}}{b'}\surd{\dfrac{\gamma^{\mathrm n}}{\delta^{\mathrm n}}}\surd{\gamma\delta},\\[7pt]
D^{\mathrm n}=\dfrac{b^{\mathrm n+1}}{b'}\surd{\dfrac{\delta^{\mathrm n}}{\gamma^{\mathrm n}}}\surd{\gamma\delta}\\
\begin{aligned}\dfrac{H}{k}y^{\frac12}\eta&=\dfrac{D}{b}\cdot\surd{\dfrac{bC}{aD}}\cdot
{}^{4}\!\surd{\dfrac{b'C'}{a'D'}}\\&\quad\cdot{}^{8}\!\surd{\dfrac{b''C''}{a''D''}}\cdot
{}^{16}\!\surd{\dfrac{b'''C'''}{a'''D'''}}.\;.\;.\end{aligned}\end{gathered}
\]
and these quantities $C,\allowbreak D$ approach the value $\frac{k}{b'}\surd{\gamma\delta}$ rapidly, while $\gamma^{\mathrm n}$ and $\delta^{\mathrm n}$ together either grow or decrease considerably as soon as $\frac{x}{k}y^{\frac12}\eta$ differs from unity.]

\clearpage
\endgroup
